% META: {"id": "TSPE-CALCUL-INTEGRAL-METH-01", "chapitre": "TSPE-CALCUL-INTEGRAL", "type_objet": "methode", "capacites_codes": ["C2"], "methodes": ["M1"], "sources_inspiration": [], "status": "approved"}
\begin{fichemethode}{M1}{Calculer une intégrale}

\quandUtiliser{%
  Utiliser cette methode pour calculer $\int_a^b f(x)\,\mathrm{d}x$.
}

\pasApas{%
  \begin{enumerate}
    \item \textbf{Chercher une primitive} $F$ de $f$ (table de reference ou forme $(v'\circ u)u'$).
    \item \textbf{Si $f$ est un produit dont aucune primitive directe n'apparait}, envisager une integration par parties : choisir $u$ (qui se simplifie en derivant) et $v'$ (dont on connait une primitive $v$).
    \item \textbf{Calculer} $F(b)-F(a)$ (ou appliquer la formule d'IPP).
    \item \textbf{Verifier} la coherence du signe du resultat avec le signe de $f$ sur $[a,b]$ (si $f\geq0$, l'integrale doit etre $\geq0$).
  \end{enumerate}
}

\end{fichemethode}
